
Results are presented in causal order: the NLI clustering behavior (RQ3) causally precedes the AUROC outcome (RQ1).

\subsubsection{5.1 NLI Clustering Aggregation (H-M2, RQ3)}

The NLI aggregation rate — the fraction of examples for which at least one pair of stochastic responses is clustered together — is \textbf{0.272 (95\% CI [0.253, 0.292])}, measured from 2,000 HaluEval-QA examples using the cluster counts produced in H-E1/H-M1. The gate threshold for H-M2 (aggregation_rate ≥ 0.50) is not satisfied; the CI upper bound (0.292) is also below the secondary threshold (0.30). Gate result: \textbf{PIVOT}.

The cluster count distribution across all 2,000 examples is:

\begin{table}[htbp]
\centering
\begin{tabular}{lll}
\toprule
Cluster Count & Frequency & Fraction \\
\midrule
1 (full collapse — all in one cluster) & 4 & 0.2\% \\
2 & 22 & 1.1\% \\
3 & 112 & 5.6\% \\
4 & 406 & 20.3\% \\
5 (no aggregation — all in distinct clusters) & 1,456 & 72.8\% \\
\bottomrule
\end{tabular}
\end{table}

Mean cluster count = 4.644 (std = 0.657, median = 5.000). In 72.8\% of examples, deberta-large-mnli classified all five stochastic responses as semantically non-equivalent, assigning each to its own cluster. The point-biserial correlation between cluster count and hallucination label was not computable (NaN) due to the near-constant cluster distribution.

\begin{figure}[htbp]
\centering
\includegraphics[width=\columnwidth]{figures/cluster_count_dist.png}
\caption{Cluster count distribution (H-M2)}
\label{fig:cluster_count_dist}
\end{figure}

\begin{figure}[htbp]
\centering
\includegraphics[width=\columnwidth]{figures/aggregation_rate.png}
\caption{NLI aggregation rate vs. gate threshold}
\label{fig:aggregation_rate}
\end{figure}

\begin{figure}[htbp]
\centering
\includegraphics[width=\columnwidth]{figures/cluster_count_cdf.png}
\caption{CDF of cluster counts}
\label{fig:cluster_count_cdf}
\end{figure}

When all N=5 responses are in distinct clusters with equal probability p = 1/5, semantic entropy equals log₂(5) ≈ 2.322 bits for every example, producing a constant signal that carries no information about hallucination status. This is the causal mechanism underlying the degenerate semantic entropy result reported in Section 5.2.

\subsubsection{5.2 Degenerate Semantic Entropy (H-M1, RQ2)}

Semantic entropy scores are constant across all 2,000 examples: std = 4.14 × 10⁻²⁵. This near-zero variance results from the 72.8\% of examples at the maximum cluster count (5/5) contributing the same entropy value (log₂(5) ≈ 2.322 bits), with only 4 examples (0.2\%) yielding cluster count = 1 (SE = 0 bits). The variation introduced by the remaining 27.2\% of examples is insufficient to produce measurable variance. The Pearson correlation r(TE, SE) is undefined (division by zero in standard deviation). Gate result: \textbf{DEGENERATE_PASS}.

The cluster distribution from H-M1 (mean clusters = 4.644) corroborates the H-M2 finding and confirms that the degenerate SE output propagates from the NLI clustering failure in H-E1.

\begin{figure}[htbp]
\centering
\includegraphics[width=\columnwidth]{figures/degenerate_summary.png}
\caption{Token entropy vs. semantic entropy scatter (degenerate SE)}
\label{fig:degenerate_summary}
\end{figure}

\begin{figure}[htbp]
\centering
\includegraphics[width=\columnwidth]{figures/cluster_count_dist.png}
\caption{Cluster count distribution (H-M1)}
\label{fig:cluster_count_dist}
\end{figure}

\subsubsection{5.3 Main AUROC Comparison (H-E1, RQ1)}

AUROC values and 95\% bootstrap confidence intervals for all three UQ methods on HaluEval-QA with LLaMA-2-7B-chat are:

\begin{table}[htbp]
\centering
\begin{tabular}{llll}
\toprule
Method & AUROC & 95\% CI Lower & 95\% CI Upper \\
\midrule
Semantic Entropy (SE) & 0.5000 & 0.5000 & 0.5000 \\
Token Entropy (TE) & 0.4829 & 0.4585 & 0.5090 \\
SelfCheckGPT-BERTScore (SCG) & 0.3562 & 0.3321 & 0.3803 \\
\bottomrule
\end{tabular}
\end{table}

\begin{figure}[htbp]
\centering
\includegraphics[width=\columnwidth]{figures/auroc_bar_chart.png}
\caption{AUROC bar chart with 95\% bootstrap CIs}
\label{fig:auroc_bar_chart}
\end{figure}

\begin{figure}[htbp]
\centering
\includegraphics[width=\columnwidth]{figures/roc_curves_overlay.png}
\caption{ROC curves overlay}
\label{fig:roc_curves_overlay}
\end{figure}

\textbf{Semantic entropy (AUROC = 0.5000, CI = [0.5000, 0.5000]).} The degenerate confidence interval — a single point at exactly 0.5 — is the expected signature of a constant signal: every threshold yields the same true positive rate and false positive rate, producing an ROC curve on the diagonal. This result is a direct consequence of the H-M2 NLI clustering failure described in Section 5.1, not a property of semantic entropy in general.

\textbf{Token entropy (AUROC = 0.4829, CI = [0.4585, 0.5090]).} Near-random. The CI spans 0.5, indicating that token entropy is consistent with zero discrimination on this benchmark and model combination.

\textbf{SelfCheckGPT-BERTScore (AUROC = 0.3562, CI = [0.3321, 0.3803]).} Below random chance. The CI upper bound (0.3803) is entirely below 0.5, indicating a statistically significant negative correlation between BERTScore consistency and HaluEval-QA hallucination labels. This is the opposite of the positive discrimination assumed by SelfCheckGPT's design.

Pairwise comparisons (Bonferroni-corrected α = 0.0167):

\begin{table}[htbp]
\centering
\begin{tabular}{lllll}
\toprule
Pair & Winner & Δ AUROC & CI Overlap & Qualifies (Δ ≥ 0.05, non-overlapping) \\
\midrule
SE vs. SCG & SE & 0.1438 & Non-overlapping & Yes \\
TE vs. SCG & TE & 0.1268 & Non-overlapping & Yes \\
SE vs. TE & SE & 0.0171 & Overlapping & No \\
\bottomrule
\end{tabular}
\end{table}

Two of three pairwise comparisons satisfy both the Δ ≥ 0.05 threshold and the non-overlapping CI criterion. Gate result: \textbf{PASS (MUST_WORK satisfied)}. However, both qualifying pairs are driven by SelfCheckGPT performing significantly below random, not by semantic entropy outperforming token entropy. The primary prediction — SE achieves ≥ 0.05 AUROC advantage over TE — is not supported (Δ = 0.0171, overlapping CIs).

\subsubsection{5.4 Label-Stratified Cluster Counts (H-M2 Supplementary)}

The cluster count by hallucination label shows similar distributions for hallucinated and factual examples, consistent with label-agnostic NLI clustering behavior.

\begin{figure}[htbp]
\centering
\includegraphics[width=\columnwidth]{figures/cluster_count_by_label.png}
\caption{Cluster count by hallucination label}
\label{fig:cluster_count_by_label}
\end{figure}

\subsubsection{5.5 Summary of Sub-Hypothesis Outcomes}

\begin{table}[htbp]
\centering
\begin{tabular}{lllll}
\toprule
Sub-Hypothesis & Type & Gate & Result & Primary Finding \\
\midrule
H-E1 & EXISTENCE & MUST_WORK & PASS & 2 qualifying method pairs (SE>SCG, TE>SCG); gap driven by SCG below-random \\
H-M1 & MECHANISM & MUST_WORK & DEGENERATE_PASS & SE std = 4.14 × 10⁻²⁵; Pearson r undefined \\
H-M2 & MECHANISM & SHOULD_WORK & PIVOT & Aggregation rate = 0.272 (CI [0.253, 0.292]); A2 violated \\
H-M3 & MECHANISM & MUST_WORK & Not executed & Mistral-7B-Instruct experiment absent \\
\bottomrule
\end{tabular}
\end{table}

